\tablehead{tab:S6}{Two-stage classifier-gated routing outcomes and
fail-closed verification, reported per device and per routing path. Paths,
devices and measurement types are never combined into a single headline
figure. Part~A is physical-device counter evidence; Part~B is emulator
behavioural-acceptance evidence and is not physical-device evidence.}

\needspace{10\baselineskip}\noindent\textbf{Part A. Physical-device harness-measured routing outcomes}

\smallskip
\noindent\textit{Xiaomi Redmi Note~14 Pro+~5G (one session while charging;
$n = 100$ per path)}

\begin{center}\footnotesize\setlength{\tabcolsep}{5pt}
\begin{tabular}{p{115pt} r r r r r}
\toprule
             &     & Det.  & Det.  & Offline & Offline \\
Routing path & $n$ & exec. & skip. & success & failure \\
\midrule
Accepted Corn        & 100 & 100 & 0   & 100 & 0 \\
Accepted Pepper      & 100 & 100 & 0   & 100 & 0 \\
Accepted Tomato      & 100 & 100 & 0   & 100 & 0 \\
OOD rejection        & 100 & 0   & 100 & 100 & 0 \\
Crop mismatch        & 100 & 0   & 100 & 100 & 0 \\
\bottomrule
\end{tabular}\end{center}

\smallskip
\noindent\textit{iPhone~15 Pro Max (three consecutive physical sessions;
$n = 100$ per run; only the out-of-distribution path was exercised in the
100-iteration harness loop in the current evidence set)}

\begin{center}\footnotesize\setlength{\tabcolsep}{5pt}
\begin{tabular}{l l r r r r r}
\toprule
    & Thermal &     & Det.  & Det.  & Offline & Offline \\
Run & state   & $n$ & exec. & skip. & success & failure \\
\midrule
Run 1 & nominal & 100 & 0 & 100 & 100 & 0 \\
Run 2 & fair    & 100 & 0 & 100 & 100 & 0 \\
Run 3 & fair    & 100 & 0 & 100 & 100 & 0 \\
\bottomrule
\end{tabular}\end{center}

{\footnotesize
Across the five Xiaomi paths and three iPhone runs, all operations completed
without an offline failure. The detector was never invoked for an
out-of-distribution image or a crop mismatch, but it was invoked for every
accepted image. These are observed counts rather than statistically estimated
failure rates.
\par}

\medskip
\needspace{10\baselineskip}\noindent\textbf{Part B. Emulator behavioural-acceptance evidence
(supplementary only; not physical-device evidence; not used to support any
deployment-performance claim)}

\begin{center}\footnotesize\setlength{\tabcolsep}{5pt}
\begin{tabular}{p{110pt} l c c l}
\toprule
     & Classifier      & Detector & Detector & Routed \\
Case & outcome         & executed & calls    & classes \\
\midrule
Supported Corn   & Corn @ 0.924    & yes & 1 & Corn diseases   \\
Supported Pepper & Pepper @ 0.891  & yes & 1 & Pepper diseases \\
Supported Tomato & Tomato @ 0.949  & yes & 1 & Tomato diseases \\
Unsupported cassava leaf & Other @ 0.800 & no  & 0 & none        \\
Crop mismatch (Corn image, Tomato selected) & Corn @ 0.924 & no & 0 & none \\
Rotated Corn image & Corn @ 0.925 & yes & 1 & Corn diseases   \\
\bottomrule
\end{tabular}\end{center}

{\footnotesize
Part~B was conducted on an Android phone emulator. Six cases were passed
through the application's normal user interface, while an independent counter
recorded whether the detector was invoked. A separate 100-iteration automated
test used a synthetic image that the classifier consistently rejected. The
detector was skipped in all 100 iterations, and every iteration completed
without failure. The manually driven cases and the automated test are
reported separately and are not combined into a single count.

\textit{Notes.} (1)~\emph{Fail-closed} means that the detector is not invoked
when the classifier rejects an unsupported image or when its prediction does
not match the crop selected by the user. In both situations, the application
still returns a valid response rather than failing. (2)~No accepted-path or
crop-mismatch-path 100-iteration harness data exists for the iPhone in the
current evidence set; this is recorded as a limitation, not filled with an
estimate. (3)~The Xiaomi evidence in Part~A is from a single pilot session
while the device was plugged in and charging, a disclosed measurement
condition. This may affect the timing figures in Table~\ref{tab:S5}, but not
the observed routing counts reported here. (4)~Part~B rows are
Android-only and were run under software emulation on a desktop host; they
are not physical-device performance evidence for the iPhone, the Xiaomi or any
other handset.
\par}
\tabletail
